\chapter{Background and Literature Review}
\label{ch:background}

\section{Breast ultrasound as an imaging modality}
Ultrasound forms images from the transmission and reception of acoustic waves. At a simplified level, the speed of a small-amplitude pressure wave in a medium is related to its bulk modulus $K$ and density $\rho$ by
\begin{equation}
c = \sqrt{\frac{K}{\rho}}.
\end{equation}
At an interface between tissues with acoustic impedances $Z_1$ and $Z_2$, the normal-incidence intensity reflection coefficient can be written as
\begin{equation}
R_I = \left(\frac{Z_2-Z_1}{Z_2+Z_1}\right)^2.
\end{equation}
Spatially varying reflection, scattering, attenuation, time-gain compensation, and scanner processing jointly determine the displayed B-mode image. The resulting pixel array is therefore not a direct tissue measurement and can vary substantially between acquisitions.

Breast ultrasound is used alongside clinical history and other imaging. Radiologists may consider features such as shape, margin, orientation, echogenicity, and posterior acoustic behaviour within the ACR BI-RADS ultrasound lexicon \citep{Mendelson2013,Stavros1995}. These descriptors provide clinical context for why geometry and local contrast matter. They do not justify inferring a BI-RADS category from this project's classifier or from its heuristic image measurements.

\section{Speckle and image variability}
Speckle arises from constructive and destructive interference among echoes from unresolved scatterers. A frequently used simplified observation model combines signal-dependent and additive components:
\begin{equation}
I(x,y)=f(x,y)\,\eta_m(x,y)+\eta_a(x,y),
\label{eq:speckle_model}
\end{equation}
where $f$ is an underlying reflectivity field, $\eta_m$ is a multiplicative term, and $\eta_a$ represents additive acquisition noise. This model is an approximation: displayed clinical images also contain beamforming, compression, filtering, and device-specific processing.

Traditional reduction methods include local-statistics filters and speckle-reducing anisotropic diffusion \citep{Lee1980,Yu2002}. More recent work evaluates learned models under perturbations to determine whether apparently strong classifiers remain stable when image quality changes \citep{Jiang2023}. Robustness must be tested on stored predictions with a specified perturbation process and seed. A visually plausible synthetic example is not by itself a performance experiment.

\section{Geometric preprocessing}
CNN batches usually require a common spatial size. If an image of height $H$ and width $W$ is directly resized to $H'$ by $W'$, its axis-specific scale factors are
\begin{equation}
s_y=\frac{H'}{H}, \qquad s_x=\frac{W'}{W}.
\end{equation}
When $s_x\neq s_y$, the transformation is anisotropic and changes apparent aspect ratios. An alternative is to crop a defined region, pad it to a square, and then resize it. If the square side is $d$, both axes use the same final factor $224/d$ in this project.

This geometric property is deterministic, but its effect on classification is empirical. Cropping can discard useful context; reflection padding can introduce repeated boundary structure; and mask-guided training can differ from unmasked deployment. A fair ablation must hold the partition, architecture, optimisation budget, and random seeds constant while changing only the preprocessing strategy.

\section{Contrast Limited Adaptive Histogram Equalisation}
Global histogram equalisation uses an image-wide cumulative distribution and can over-emphasise background variation. Contrast Limited Adaptive Histogram Equalisation (CLAHE) instead applies local mappings and clips histogram mass before redistribution \citep{Zuiderveld1994}. For a contextual region with histogram $h(k)$, a conceptual clipped histogram is
\begin{equation}
h_c(k)=\min(h(k),\beta)+r(k),
\end{equation}
where $\beta$ is the clip threshold and $r(k)$ distributes the clipped mass. Interpolation between neighbouring mappings reduces tile-boundary discontinuities.

The implementation in this repository applies OpenCV CLAHE to the luminance channel in LAB colour space, with a clip limit of 2.0 and a $16\times16$ tile grid. This is a documented configuration choice, not evidence that CLAHE improves discrimination. Its effect requires a controlled ablation and should include both predictive and image-quality analysis.

\section{Convolutional representation learning}
CNNs learn filters through repeated convolution, non-linearity, and aggregation. For an activation tensor $\mathbf{X}^{(l-1)}$, a layer can be represented schematically as
\begin{equation}
\mathbf{X}^{(l)}=\sigma\!\left(\mathbf{W}^{(l)}*\mathbf{X}^{(l-1)}+\mathbf{b}^{(l)}\right).
\end{equation}
Medical-image datasets are often much smaller than general computer-vision corpora, so transfer learning is commonly used to initialise useful low- and mid-level features \citep{Litjens2017}.

\subsection{ResNet-50}
Residual networks learn a residual function around an identity path,
\begin{equation}
\mathbf{y}=\mathcal{F}(\mathbf{x},\{\mathbf{W}_i\})+\mathbf{x},
\end{equation}
which supports optimisation of deeper architectures \citep{He2016}. In this project, the original classification layer is replaced by dropout and a two-output linear layer.

\subsection{EfficientNet-B0}
EfficientNet uses coordinated scaling of network depth, width, and input resolution, with the B0 model serving as the baseline architecture \citep{Tan2019}. Its mobile inverted bottleneck blocks include squeeze-and-excitation operations that reweight channels \citep{Hu2018}. The project replaces the classifier with dropout and a two-output linear layer.

\subsection{Custom CNN baseline}
The custom model provides a compact, from-scratch baseline with four convolutional stages, batch normalisation, ReLU activations, max pooling, adaptive average pooling, dropout, and a two-output layer. It helps distinguish the contribution of transfer learning from that of the surrounding pipeline, but a fair comparison requires equivalent split integrity and a predeclared tuning budget.

\section{Outputs, confidence, and explanation}
The two logits are converted to a mutually exclusive probability vector using softmax:
\begin{equation}
p(y=c\mid\mathbf{x})=\frac{\exp(z_c)}{\sum_{j=1}^{K}\exp(z_j)}.
\end{equation}
Softmax normalisation does not guarantee calibration. Neural networks can be overconfident, so a displayed score must not be interpreted as diagnostic certainty \citep{Guo2017}. Calibration should be evaluated on an independent validation design if the model is developed further.

Grad-CAM uses gradients flowing into a convolutional layer to form a coarse localisation map \citep{Selvaraju2017}. Such maps can support debugging and hypothesis generation, but they neither prove causal reasoning nor validate clinical relevance. This project labels them as research visualisations.

\section{Dataset provenance and evaluation design}
The BUSI dataset paper describes 780 images with normal, benign, and malignant labels and associated masks \citep{AlDhabyani2020}. OASBUD includes multiple views and radio-frequency-derived information for breast lesions \citep{Piotrzkowska2017}. These characteristics make group-level separation essential. When related views cross partitions, a model may encounter subject-specific appearance during both optimisation and evaluation.

Reproducible evaluation therefore requires more than a random seed. A defensible record includes original subject identifiers, grouping rules, dataset version and licence, checksums, exclusions, split manifest, preprocessing configuration, model configuration, training seed, selected checkpoint, and per-image predictions. Only then can summary metrics be independently regenerated.

\section{Literature synthesis and research gap}
The literature establishes several relevant components: public breast-ultrasound datasets \citep{AlDhabyani2020,Piotrzkowska2017}; deep representation learning for medical images \citep{Litjens2017,Cheng2016}; efficient and residual transfer-learning architectures \citep{He2016,Tan2019}; contrast enhancement \citep{Zuiderveld1994}; and perturbation-based robustness testing \citep{Jiang2023}. The gap addressed here is practical integration with an explicit evidence boundary. A classifier, dashboard, and polished report can appear complete while still lacking the provenance needed to justify their headline numbers.

Accordingly, this dissertation treats data auditability as part of the modelling method. It implements the full software path, reports the current executable output as provisional, and defines the group-aware experiment required before comparative or clinical conclusions can be drawn.
